\documentclass[11pt]{amsart}
\usepackage{amsmath,amssymb,mathrsfs,hyperref}
\newtheorem{theorem}{Theorem}[section]
\newtheorem{proposition}[theorem]{Proposition}
\newtheorem{lemma}[theorem]{Lemma}
\newtheorem{remark}[theorem]{Remark}
\newcommand{\Fp}{\mathbb F_p}
\newcommand{\Hom}{\operatorname{Hom}}
\newcommand{\Tr}{\operatorname{Tr}}
\newcommand{\Sp}{\operatorname{Sp}}
\newcommand{\SU}{\operatorname{SU}}
\newcommand{\SL}{\operatorname{SL}}
\newcommand{\GL}{\operatorname{GL}}
\newcommand{\U}{\operatorname{U}}
\newcommand{\Heis}{\operatorname{Heis}}
\newcommand{\FPdim}{\operatorname{FPdim}}
\title{Exact flux braid images in degree-three permutation gaugings}
\author{Titania}
\date{14 November 2026}
\begin{document}
\begin{abstract}
For every prime $p\ge5$, we determine the exact braid images on all
fusion spaces of identical three-cycle fluxes in every $C_3$ permutation
gauging of a pointed modular category on $\mathbb F_p$.
We identify the actual equivariant charge projectors and crossing
scalars. Generic closed sectors have determinant-$\mu_6$ linear or
unitary images modulo the order-three deck group; charged sectors have
full symplectic or Lagrangian translation groups with explicit
Heisenberg--Weil extensions. We give projective and linear braid-kernel
criteria, restore every diagonal flux species, compare all three gauging
cocycles, and determine the small-strand exceptions.
\end{abstract}
\maketitle

\section{Statement and notation}
Fix a prime $p\ge5$, $p\ne3$, and a metric form
$q(x)=\exp(2\pi i c x^2/p)$ on $\Fp$, $c\ne0$.
Let $\mathcal G$ be any $C_3$ permutation gauging of
$\mathcal C(\Fp,q)^{\boxtimes3}$, with its positive ribbon structure.
All three gauging choices are included. A three-cycle flux has grade
$a=1$ or $2$ relative to a fixed generator of $C_3$.
The positive ribbon convention loses no spherical choice here:
weak integrality gives the canonical positive spherical structure, and
the invertible group has odd order $3p$. Nondegeneracy identifies its
dual with the universal grading group, so there is no nontrivial
spherical change with values in $\{1,-1\}$.

Put $E=\Fp[z]/(z^2+z+1)$, with $\bar z=z^{-1}$.
It is $\mathbb F_{p^2}$ for $p\equiv2\pmod3$, and $\Fp\oplus\Fp$
otherwise. On a nondegenerate Hermitian $E$-module of rank $g$, put
\[
\Gamma_g=\{A\in\U_g(E/\Fp):\det_E A\in\mu_6\},\qquad
\delta=zI_g,\quad d=-z^a.
\]
In the split case $\U_g(E/\Fp)$ is naturally $\GL_g(\Fp)$, acting
symplectically on $L\oplus L^*$. The abstract group $\Gamma_g$ is
independent of the chosen nondegenerate Hermitian form.
Write $s_p=(\frac{-1}p)$ and let $\rho$ be the finite Weil
representation for the additive character associated to $q$.
The choice of symplectic basis does not change any group below.

There is a canonical pointed factor and a reduced gauging
\begin{equation}\label{split}
\mathcal G\simeq\mathcal C(\Fp,q^3)\boxtimes\mathcal E,
\qquad\FPdim(\mathcal E)=9p^2.
\end{equation}
The neutral metric plane of $\mathcal E$ is
$W=\{(x_1,x_2,x_3):\sum x_i=0\}$, with cyclic rotation.
Its electric objects are $e^j$, $j\in\mathbb Z/3$.
Its other neutral simples are the three-element orbits $O_b$,
$0\ne b\in W$. There are three flux simples in each nonneutral grade,
all of dimension $p$. Write a flux in $\mathcal G$ as $X=x\boxtimes Z$,
where $Z\in\mathcal E$ has grade $a$ and twist $\theta$.
For $Y=y\boxtimes U$, a nonzero fusion space forces $y=nx$.
The diagonal factor multiplies each braid generator by $q(x)^3$.
We state the reduced answer first and then specify the resulting groups
and kernels for $\mathcal G$.

For $n\ge2$ set
\[
g=\begin{cases}n-2,&3\mid n,\\n-1,&3\nmid n.\end{cases}
\]
Let $A_n:B_n\to\Gamma_g$ be the compact Burau representation at
$z^a$, and let $w:B_n\to\mathbb Z$ send every standard crossing to one.
The matrices and their compact quotient are specified in Section~\ref{arith}.

\begin{theorem}\label{main}
For $n\ge2$ the following describe every nonzero reduced fusion space
$\Hom_{\mathcal E}(U,Z^n)$.
\begin{enumerate}
\item Suppose $3\nmid n$. Then $U$ is a flux of grade
$r=an\bmod3$, represented by $r\in\{1,2\}$. Put
\[
\chi=(\theta_U/\theta^{n^2})^r\in\mu_3.
\]
The fusion space is the $\chi$-eigenspace of $\rho(\delta)$, and the
braid action is
\begin{equation}\label{closedaction}
\beta\longmapsto (s_p\theta)^{w(\beta)}
\rho(A_n(\beta))\big|_{\mathscr W_\chi}.
\end{equation}
\item Suppose $3\mid n$ and $U=e^j$. The same assertion holds with
\[
\chi=e^{2\pi i j/3}\theta^{an}.
\]
In both cases $\theta^9=1$. For $g\ge2$ the projective image is exactly
$\Gamma_g/\langle\delta\rangle$, with projective kernel
\[
\{\beta:A_n(\beta)\in\langle\delta\rangle\}.
\]
The dimension is $(p^g+2\varepsilon^g)/3$ for $\chi=1$, and
$(p^g-\varepsilon^g)/3$ otherwise, where $\varepsilon=(\frac{-3}p)$.
\item Suppose $3\mid n$ and $U=O_b$, $b\ne0$. The dimension is $p^g$.
There is a degenerate Burau representation
\[
\widetilde A_n:B_n\longrightarrow E^g\rtimes\Gamma_g.
\]
Its translation part, paired with the boundary character of $b$, has
image $K_b$. Here $K_b$ is the full symplectic additive space $\Fp^{2g}$
if $E$ is a field or $b$ is nonisotropic; in the remaining split
isotropic case it is a Lagrangian $\Fp^g$.
The projective image is exactly $K_b\rtimes\Gamma_g$.
Its projective kernel consists precisely of braids with compact part
$I$ and radical translation annihilated by the boundary character.
\end{enumerate}
For $g=1$, the compact group is $\mu_6$ and the generic projective
image in the first two cases is $C_2$, with kernel $w^{-1}(2\mathbb Z)$.
Its linear image is cyclic of order
$2\operatorname{ord}(\theta\chi^a)$. The only exception is
$p=5$, $\chi=1$: the fusion space is one-dimensional and its linear
image is $\langle\theta\rangle$, so its projective image is trivial.
Its linear kernel is $w^{-1}(\operatorname{ord}(\theta)\mathbb Z)$.
For $n=0,1$ all nonzero spaces are one-dimensional and the braid group
is trivial. Scalar extensions in every case are specified below.
\end{theorem}

Here is an explicit finite-group description of the linear extensions;
it also makes their scalar kernels computable without an implicit
choice of braid character. Put $\lambda=s_p\theta$ and
\begin{equation}\label{hat}
\widehat\Gamma_{g,a}=
\{(A,k)\in\Gamma_g\times\mathbb Z/18:\det_E A=d^k\}.
\end{equation}
Its order is $3|\Gamma_g|$.
For a closed-sector character $\chi$ in the generic range of
Theorem~\ref{main}, the linear image is
\begin{equation}\label{closedextension}
\widehat\Gamma_{g,a}/
\{(\delta^j,k):\det_E\delta^j=d^k,\ \lambda^k\chi^j=1\}.
\end{equation}
The linear braid kernel is given by $A_n(\beta)=\delta^j$ and
$\lambda^{w(\beta)}\chi^j=1$.
For charged sectors define
\[
N_\theta=\{(I,k):k\equiv0\pmod6,\ \lambda^k=1\},\qquad
H_\theta=\widehat\Gamma_{g,a}/N_\theta.
\]
It has order $|\Gamma_g|$ if $\theta^3=1$, and $3|\Gamma_g|$ otherwise.
The full-translation charged image is
\begin{equation}\label{chargedextension}
\Heis_p(\Fp^{2g})\rtimes H_\theta,
\end{equation}
in its faithful Schr\"odinger--Weil realization. Its scalar subgroup is
$\mu_p$, enlarged by $\mu_3$ exactly when $\theta^3\ne1$.
The split isotropic image is $\Fp^g\rtimes H_\theta$, with scalar
subgroup trivial or $\mu_3$ according to the same condition.
The action of $H_\theta$ on the Heisenberg group or translation space
is through its projection to $\Gamma_g$.

\begin{proposition}[Restoring every original flux species]\label{species}
Assume $n\ge2$.
For $X=x\boxtimes Z$ put $\kappa=q(x)^3$. If $x=0$ the preceding
groups and kernels already give the answer. If $x\ne0$, then $\kappa$
has order $p$. Projective images and kernels are unchanged. For every
closed sector, including the one-dimensional exception, the linear image
is the direct product of $\mu_p$ with its reduced image. The same is
true for a split isotropic charged sector. Its linear kernel is the
intersection of the reduced linear kernel with $w^{-1}(p\mathbb Z)$.
For a full-translation charged sector the linear image is the same
matrix group as the reduced image. In its word criterion replace
$\lambda^w\psi(z)$ by $\kappa^w\lambda^w\psi(z)$.
\end{proposition}
\begin{proof}
For a reduced crossing its sixth power is $\lambda^6$, so the
eighteenth power is one. The original crossing has eighteenth power
$\kappa^{18}$, which generates $\mu_p$. This isolates the extra scalar
and recovers the reduced generators whenever $n\ge2$. In the closed
and isotropic cases their scalar subgroups have order prime to $p$,
so the product is direct and the kernel condition follows. In the
full-translation case $\mu_p$ is already contained in the reduced
image, and hence multiplying generators by its elements changes
neither the generated matrix group nor its abstract extension.
The braid groups for $n=0,1$ are trivial.
\end{proof}

\begin{remark}[The first six-strand vacuum and the prime-five control]
For grade one, the six-strand vacuum is the
$\theta^6$-eigenspace in the genus-four Weil module. Its actual
projector is
\[
\frac13\sum_{k=0}^2\theta^{3k}\rho(\delta)^k.
\]
At $p=5$, its dimension is $209$ if $\theta^3=1$, and $208$
otherwise; its projective image is $\U_4(5)/\mu_3$ in all cases.
Its reduced linear image is $\U_4(5)/\mu_3$ if $\theta=1$,
is $\U_4(5)$ if $\theta\ne1$ but $\theta^3=1$, and has order
$3|\U_4(5)|$ and scalar subgroup $\mu_9$ otherwise, with the
explicit extension in~\eqref{closedextension}.

As a consistency check, the stored dimension-$225$, rank-$17$
fusion tensors in the Anyon Atlas~\cite{Atlas} give vacuum multiplicities
$(2,208)$ at strand counts $(3,6)$ for B17-00346 and B17-00347,
whose flux twists have nontrivial cube, and $(1,209)$ for B17-00506,
whose flux twist cubes are one. This checks the projector's multiplicity
prediction, not the categorical identification or exact braid image.
The catalogue's categorification flags for these entries are false;
the construction of the anisotropic gauging itself is prior~\cite{GPR}.
\end{remark}

\section{Relation to previous work}
We use previously established existence and finiteness results. The
cover modular functor is due to Barmeier--Schweigert~\cite{BS}; its
weak rigidity is upgraded by Deshpande--Mukhopadhyay~\cite{DM},
Corollary 2.11. Crossed extension/zesting theory supplies the three
gauging choices~\cite{DGP RZ}. In the anisotropic case the reduced
categories themselves, their fusion rules and their non-group-theoreticality
are the Galindo--Plavnik--Rowell construction~\cite{GPR}.
The pointed closed-surface representation and the finite Weil
linearization are prior~\cite{Stirling,GH}.
Barkeshli--Jian--Qi~\cite{BJQ}, Section VI, equations (91)--(98),
already construct the order-three genon Wilson algebra and its braid
automorphisms, including the three-defect torus interpretation in
Section VI.B. Their Section VII computes the gauged order-two projection
and proposes the analogous order-three relation. We build on this
geometric and Wilson-algebra method, retaining the actual order-three
equivariant projectors, peripheral characters, crossing phases and exact
image and kernel information.
Exact finite-coefficient
trielliptic monodromy was studied by Achter--Pries~\cite{AP}; its raw
stack versus determinant-normalization convention is not being used
as an unqualified linear input here. The explicit Burau model is
prior~\cite{Ven}. We use the deep classification of irreducible
transvection groups~\cite{ZS}, rather than claim that classification.
The advance is the actual charge projection, uniform projected
kernel, scalar normalization and charged affine image for this entire
family. The general fixed-subalgebra principle in~\cite{GKRZ},
Principle 5.2, is a proposal, not a prior theorem supplying these bridges.
The newer metric-tower finiteness result~\cite{GR26}, Proposition 8.2
and Theorem 8.4, gives a finite Clifford extension with its translation
and linear-image subgroups defined by intersection and image. It does
not identify the actual gauged charge spaces or the exact groups here.

\section{The categorical spaces and their Wilson algebra}\label{cat}
\begin{lemma}\label{diagonal}
The decomposition~\eqref{split} exists. The reduced neutral category
has the stated metric plane, and a nontrivial grade has one underlying
defect $T_a$ of dimension $p$. Its products are
\[
T_aT_b=pT_{a+b}\ (a+b\ne0),\qquad
T_aT_{-a}=\bigoplus_{v\in W}v.
\]
\end{lemma}
\begin{proof}
Neutral equivariant invertibles form an abelian group extension of
the diagonal $\Fp$ by $\widehat C_3$. There are three lifts of each
fixed invertible since $H^2(C_3,\mathbb C^\times)=0$.
The unique Sylow $p$ subgroup has twist $q^3$, hence is nondegenerate.
Its M\"uger centralizer gives~\eqref{split}. The electric subcategory
centralizes it, and de-equivariantization identifies the neutral
centralizer with its orthogonal complement in $\Fp^3$, namely $W$.
Rotation is fixed-point-free on $W$.
The original geometric permutation category has $p$ simples in each
nonidentity grade, since its sector is one copy of
$\mathcal C(\Fp,q)$ as an abelian category. De-equivariantizing
the product decomposition by the electric $C_3$ leaves the pointed
factor unchanged. Ranks multiply, so the reduced nonidentity grade
has exactly one simple. All homogeneous components have dimension
$p^2$, so its dimension is $p$. Dimension and grading
give the displayed products; every neutral invertible occurs once in
$T_aT_a^*$ by Frobenius reciprocity.
\end{proof}

The cover construction in~\cite{BS}, Theorem 26, assigns the pointed
modular functor of the associated cover to a marked $C_3$-cover.
Its duality represents $\Hom(1,T\otimes-)$, the same r-category
duality used in~\cite{DM}. Corollary 2.11 there makes this finite
semisimple monoidal category rigid because its neutral component is
rigid. Reconstruction identifies its genus-zero braiding and the
action on multiplicities with the lifted cover diffeomorphisms, including
simultaneous deck changes of the boundary markings. Thus equivariantization
selects invariants of the actual tensor-induced action, not the whole
cover state space. All other gaugings differ by the crossed cocycles
treated in Section~\ref{zesting}.

The Riemann--Hurwitz formula gives the genus in Theorem~\ref{main}.
For $3\nmid n$ the outgoing flux contributes one ramified cap; for
$3\mid n$ there are three unramified caps. An electric output caps
them by units. A neutral label $b=(b_1,b_2,b_3)$ caps them by the
corresponding pointed labels, whose sum must vanish.

\begin{lemma}[Wilson quantization]\label{wilson}
On a genus-$g$ pointed modular-functor space with three caps labeled
$b_1,b_2,b_3\in\Fp$, $\sum b_i=0$, the normalized Wilson operators
form the twisted homology algebra
\[
W_hW_k=\psi(\tfrac12\omega(h,k))W_{h+k},\qquad
\psi(t)=\exp(2\pi i\,2ct/p),
\]
where $h,k\in H_1(\Sigma\setminus\{P_1,P_2,P_3\},\Fp)$ and
$\omega$ is intersection. Its radical $R$ has dimension two.
A peripheral class $r=(r_1,r_2,r_3)$, modulo the diagonal, acts by
\begin{equation}\label{boundarychar}
\chi_b(r)=\psi\left(\sum r_i b_i\right).
\end{equation}
After this central specialization the algebra is
$\operatorname{End}(\mathbb C[\Fp^g])$. Lifted diffeomorphisms act
on it by their homology action. The analogous closed space is the
finite Heisenberg module with projective finite Weil action.
\end{lemma}
\begin{proof}
Use the strict pointed model with trivial associator. Its braiding is
$\exp(2\pi i c xy/p)$ and its twist is $q(x)$.
Cut a handlebody to the usual
handle graph and an arc joining the charged caps. Its basis is labeled
by $g$ independent elements of $\Fp$, so the state-space dimension
is $p^g$. Wilson curves along the handle meridians multiply a basis
vector by its double-braiding character, and curves along the
longitudes shift its label. Normal ordering with the half-intersection
factor gives exactly the displayed multiplication law. The same
skein calculation resolves any crossing of two colored curves; fusion
adds their labels. A contractible curve acts by the dimension one,
and moving a peripheral curve past a cap gives its double braiding,
which is~\eqref{boundarychar}. The relation between the three
peripheral classes is their sum, consistent with $\sum b_i=0$.
These relations show that the normalized operator depends only on
the indicated mod-$p$ homology class. Handle shifts and characters
are the standard Weyl basis of the full matrix algebra, proving
irreducibility and the stated specialization.

Naturality of insertion of a ribbon curve in a mapping cylinder sends
$W_h$ to $W_{f_*h}$; blackboard framing and the half-intersection
normalization are preserved. Schur's lemma therefore determines the
mapping-cylinder operator up to one scalar. This proves the
Heisenberg covariance assertion, including charged caps. In the
closed case it is also the projective equivalence in~\cite{Stirling},
Theorem 6.79, followed by the finite Stone--von Neumann construction
in~\cite{GH}, Section 1. No assertion that charged surfaces factor
through closed homology is made.
\end{proof}

For a nonzero orbit output, the three summands with labels
$b,zb,z^2b$ are cyclically permuted by equivariance. Invariants of
their direct sum identify with one summand. Braids commute with this
equivariance and preserve the three marked outgoing sheets, so the
gauged braid action is exactly its action on this one charged space.

\section{Finite-field Burau images}\label{arith}
Let $t=z^a$ and $m=n-1$. On the chain $E^m$ use the Hermitian matrix
with diagonal $1$, upper off-diagonal $t^{-1}$, lower off-diagonal $t$.
The standard generators are
\begin{equation}\label{braidmatrices}
r_i\epsilon_i=-t\epsilon_i,\quad
r_i\epsilon_{i-1}=\epsilon_{i-1}+t\epsilon_i,\quad
r_i\epsilon_{i+1}=\epsilon_i+\epsilon_{i+1},
\end{equation}
fixing other basis vectors. These are the specialized Burau matrices;
the cover-homology identification follows from the explicit Gassner
model in~\cite{Ven}, Lemma 14 and Section 7. The determinant recurrence
$D_m=D_{m-1}-D_{m-2}$, $D_0=D_1=1$, shows a one-dimensional radical
exactly when $m\equiv2\pmod3$. It is generated by
\[
v=\sum_{i=1}^m(1-t^i)\epsilon_i.
\]
It is fixed pointwise. The compact homology is the quotient by this
radical. The full twist acts on a nondegenerate chain for $n$ strands
by $t^n I$, hence as $\delta^{an}$.

\begin{lemma}\label{finiteimage}
For $n\ge4$ the compact image is exactly $\Gamma_g$ and its pure
image has determinant in $\mu_3$ and contains the full special group.
For $n=2,3$ the compact image is $\mu_6$. If $3\mid n$, the
uncapped degenerate image is exactly $E^g\rtimes\Gamma_g$.
\end{lemma}
\begin{proof}
Here is a direct finite-field argument, using the
PRIOR transvection classification in~\cite{ZS}.
Set $u=r_1r_2^{-1}$. Direct multiplication on three consecutive
vectors gives
\[
[u,r_1ur_1^{-1}]-I
=-(2+t)(-t,1,0)^t(0,0,1).
\]
This is a nonzero square-zero rank-one matrix: $N(2+t)=3$.
It fixes all later basis vectors. The braid commutator is pure,
since the permutations of $u$ and its $r_1$-conjugate are opposite
three-cycles. Thus a transvection belongs to the pure image, over
each field component of $E$.

The nondegenerate chain image is absolutely irreducible. An invariant
subspace containing a nonfixed vector contains an image of some
$r_i-I$, hence $\epsilon_i$, and neighboring reflections then give
all basis vectors. Nondegeneracy excludes a common fixed vector.
It is primitive: if a rank-one reflection permutes direct-sum blocks
nontrivially, the rank bound forces an interchange of two lines.
Its eigenvalues there must be $1,-1$. The exceptional eigenvalue
here is $-t\ne-1$, a contradiction.

Let $H$ be the normal closure of the displayed transvection in the
full chain image. Clifford theory makes the restriction to $H$
semisimple; primitivity permits just one isotypic component, and
the rank-one action of a transvection forces multiplicity one.
Thus $H$ is irreducible. The main theorem of~\cite{ZS} leaves natural
special-linear, symplectic and special-unitary groups over subfields.
In the split case the ground field is prime. The symplectic option
cannot have these reflections in its normalizer: its alternating
form is unique up to scalar, while in dimension at least four a
semisimple reflection with exceptional eigenvalue different from one
cannot preserve a nondegenerate alternating form up to scalar.
Indeed its fixed hyperplane has nonzero restricted form, forcing the
similitude multiplier to be one, and then its exceptional line is
orthogonal to the entire space.

In the inert case the only fields are $\Fp$ and $\mathbb F_{p^2}$.
Special-linear groups on their natural module have no invariant
Hermitian form here: over $\Fp$ its coefficients would be invariant
bilinear forms, and over $\mathbb F_{p^2}$ the unrestricted root-group
parameter violates the unitary transvection condition. The symplectic
possibilities are excluded by the same normalizer argument. Therefore
$H$ is the full special-unitary group. Since the original transvection
was pure, so is its whole normal closure. Determinants of the
$r_i$ generate $\mu_6$, and those of $r_i^2$ generate $\mu_3$,
which proves the asserted exact images.

When the original chain is degenerate, the first $g=m-1$ basis
vectors span a nondegenerate complement to its radical. Its braid
subgroup already gives the full compact group just proved. The
rank-one compact cases follow directly from~\eqref{braidmatrices}.

For the final assertion write the degenerate module as $R\oplus V$
using this complement. Its isometries fixing $R$ pointwise have form
$E^g\rtimes\U(V)$. The first $n-1$ strands supply the block-diagonal
$\Gamma_g$. Their full twist is $1$ on $R$ and $t^{-1}$ on $V$.
The last reflection has a radical row $\ell$ nonzero in both split
components: the last coefficient $1-t^m=1-t^{-1}$ of the radical
vector is a unit, and its action on $\epsilon_{m-1}$ adds
$t\epsilon_m$. Commuting it with the preceding full twist therefore
produces a nonzero radical translation in both components.

In the inert case the natural $\Gamma_g$-module remains irreducible
over $\Fp$. For $g\ge3$ it is absolutely irreducible over $E$ and
is not equivalent to its Frobenius conjugate: the special-unitary
diagonal element with entries $d,d,d^{-2},1,\ldots$ has trace not in
$\Fp$, since $d+d^{-1}=1$. For $g=1$ use $d$ itself. The standard
restriction-of-scalars argument now gives irreducibility. In the
split case the two radical modules are the natural and dual modules;
they are irreducible and nonisomorphic for $g\ge3$, and their
one-dimensional weights differ for $g=1$. A stable additive subgroup
with nonzero projections to both therefore contains both. Conjugating
the radical translation thus gives all of $E^g$, proving the split
semidirect product. For $n=3$ the same conclusion follows directly
from $r_1r_2^{-1}=I+(-t,1)^t(1,t)$ and conjugation by $r_1$.
\end{proof}

\section{Ribbon trace and the actual deck projector}
\begin{lemma}\label{normalization}
On every relevant raw cover space the crossing is the trace-normalized
Weil or affine-Weil implementer times $\lambda=s_p\theta$.
Moreover $\theta^9=1$, and the actual projectors are exactly those
in Theorem~\ref{main}.
\end{lemma}
\begin{proof}
The two-strand multiplicity space $M=\Hom(T_{2a},T_a^2)$ has
dimension $p$ and is the direct sum of the three gauged flux-output
spaces. All three output objects have dimension $p$. The ribbon
trace of the crossing is the positive-curl invariant $p\theta$.
Thus the ordinary trace of its operator on $M$ is $\theta$.
Its compact half-braid has characteristic polynomial $x^2-x+1$,
so the standard Weil character formula gives trace $s_p$.
The scalar is therefore $s_p\theta$. Factoring the first two defects
from any larger raw multiplicity space multiplies both traces by
$p^{g-1}$ and gives the same scalar.

For charged caps the homology reflection induces an affine map with
a fixed point. Indeed its radical row is supported on the moving
line of the same reflection, so its symplectic translation lies in
$\operatorname{im}(I-A)$. Its implementer is conjugate by a Weyl
operator to $\rho(A)$, has order six, and has trace $s_pp^{g-1}$.
The identical categorical trace fixes its scalar too.

Suppose $3\nmid m$, and choose an equivariant flux output $U$ of
grade $r=am\bmod3$. On the raw space let $Q$ be the honest order-three
equivariance operator relative to $U$. A $\zeta^j$-eigenvector is a
morphism from $Ue^j$. Its twist is $\theta_U\zeta^{rj}$, because
gauged twist uses equivariance at the object's grade. Thus balancing
identifies the full braid twist on the raw space with
$\theta_U\theta^{-m}Q^r$.
On the other hand the normalized braid formula gives
$\theta^{m(m-1)}\rho(\delta)^r$; the sign disappears since
$m(m-1)$ is even. Therefore
\[
Q^r=(\theta^{m^2}/\theta_U)\rho(\delta)^r.
\]
Taking $m=2$, $U=Z^*$ and cubing forces $\theta^9=1$.
The same identity for general $m$ shows its scalar is in $\mu_3$;
raising to the inverse exponent $r$ yields the first stated projector.

For $3\mid n$, pivotal currying is equivariant and identifies
$\Hom(e^j,Z^n)$ with $\Hom(Z^*e^j,Z^{n-1})$. The underlying capped
covers agree, with the same deck generator. Apply the preceding formula
using $\theta_{Z^*e^j}=\theta\zeta^{-aj}$ and $r=-a\bmod3$.
Reducing exponents modulo nine gives
$Q=\zeta^{-j}\theta^{-an}\rho(\delta)$. Its invariant space is
precisely the stated $\chi$-eigenspace.
\end{proof}

\section{Projected kernels and scalar extensions}
\begin{lemma}\label{projectedkernel}
For $g\ge2$, the projective kernel of the centralizer of $\delta$
on every Weil deck eigenspace is $\langle\delta\rangle$.
\end{lemma}
\begin{proof}
The Weyl basis gives
$|\Tr\rho(A)|^2=p^{\dim\ker(A-I)}$ by taking the trace of conjugation.
The fixed-point-free Weil character formula~\cite{GH}, equation (1.2),
gives $\Tr\rho(\delta)=\Tr\rho(\delta^2)=\varepsilon^g$, yielding
the dimensions in the theorem. If $A$ centralizes $\delta$ and is
not a deck scalar, each $A\delta^j-I$ has fixed space of dimension
at most $2g-2$. In the inert case it is a proper $E$-space. In the
split case the dimensions in the dual eigenspaces agree and their
sum is even. Averaging the three traces on the projector gives
\[
|\Tr(\rho(A)|_{\mathscr W_\chi})|\le p^{g-1}
<(p^g-2)/3\le\dim\mathscr W_\chi.
\]
A scalar unitary action would have trace modulus equal to the
dimension. Deck scalars do act scalarly, proving the assertion.
\end{proof}

The homological image and Lemmas~\ref{normalization}--\ref{projectedkernel}
give the first two image and kernel statements. Every $r_i^6=I$
and $w(\sigma_i^6)=6$, so the map to~\eqref{hat} is surjective.
An element of this group acts as $\lambda^k\rho(A)$ on a selected
eigenspace. Lemma~\ref{projectedkernel} identifies its kernel exactly
as~\eqref{closedextension}.

For $g=1$, $r^2=\delta^{2a}$, so on $\mathscr W_\chi$ its eigenvalues
belong to $\{\chi^a,-\chi^a\}$. Both signs occur unless the space
is one-dimensional: the same trace bound is at most one, whereas
every eigenspace has dimension at least two for $p\ge7$.
For $p=5$ the dimensions are $1,2,2$; on the unique one-dimensional
space the finite Weil operator $r$ acts by one. This follows either
from its character or the inert-torus multiplicity formula in
\cite{GH}, Theorem 1.3: among the six torus characters exactly the
quadratic one is omitted. This proves all small-range assertions.

\section{Charged affine groups and their scalars}
In the splitting $H_1=R\oplus V$, a braid acts as
$(r,v)\mapsto(r+\ell(v),Av)$. Lemma~\ref{wilson} specializes this to
\[
W_v\longmapsto\chi_b(\ell(v))W_{Av}.
\]
Its projective implementer is unique. It is scalar if and only if
$A=I$ and $\chi_b\circ\ell=0$, since the Weyl basis is linearly
independent. The finite image in Lemma~\ref{finiteimage} now gives
the claimed projective groups. If $E$ is a field, any nonzero
$E$-linear functional is onto $E$, so composition with the nontrivial
boundary character is injective on the translation radical. In the
split case its two components survive exactly when both components
of the output label survive. This is the stated distinction between
isotropic and nonisotropic outputs.

For explicit linear kernels use the Heisenberg convention
\[
\pi(v)\pi(v')=\psi(\tfrac12\omega(v,v'))\pi(v+v').
\]
For each braid generator choose $x_i$ with
$t_i=(I-A_i)x_i$ its specialized translation. Its canonical affine
implementer is
\[
C_i=\pi(x_i)\rho(A_i)\pi(x_i)^{-1}
=\psi(z_i)\pi(t_i)\rho(A_i),\quad
z_i=-\tfrac12\omega(x_i,A_ix_i).
\]
The phase $z_i$ is independent of the fixed-vector ambiguity in $x_i$.
Words are evaluated in the finite group law
\[
(t,A,z)(t',A',z')=
(t+At',AA',z+z'+\tfrac12\omega(t,At')).
\]
The actual operator is $\lambda^{w(\beta)}\psi(z(\beta))
\pi(t(\beta))\rho(A(\beta))$. Consequently its linear kernel is
$A=I$, $t=0$ and $\lambda^w\psi(z)=1$. The two scalar orders are
coprime, so the last condition is equivalently $\lambda^w=1$, $z=0$.
This gives a finite-field word criterion in addition to the projective
kernel in Theorem~\ref{main}.
As a scalar consistency check, a full braid twist in this charged
sector has $A=I$, $t=0$. Balancing forces its remaining Heisenberg
phase to be the output twist $q_W(b)$, since $9\mid n^2$.
It is trivial on isotropic outputs and generates $\mu_p$ on every
nonisotropic output.

For full translations, commutators of translation lifts yield the
entire central $\mu_p$. Their non-$p$ scalar parts have order at most
three since their compact determinant is one. Taking cube powers
therefore isolates the full Heisenberg $p$ subgroup. Removing its
translation and phase from each crossing leaves $\lambda\rho(A_i)$.
This complement is exactly $H_\theta$: its compact part is faithful
projectively on the full Heisenberg module and its scalar kernel
over $I$ is $N_\theta$. Its intersection with the Heisenberg group
is trivial by the coprime scalar orders. This proves~\eqref{chargedextension}.

In the split isotropic case choose a Schr\"odinger polarization for
which the surviving translations are shifts. The compact group
preserves this Lagrangian, and its Weil operator is a permutation
times the quadratic determinant character. Each affine reflection
has a fixed point in this same Lagrangian, so conjugating by a shift
introduces no scalar phase. Trace normalization makes every actual
crossing a permutation times $\theta$. Thus no $\mu_p$ scalar is
generated. The same cube-power isolation of translations and the
same determinant fiber product give $\Fp^g\rtimes H_\theta$.

\section{Cocycle comparison}\label{zesting}
For completeness the intrinsic twist formulas can be related to an
arbitrary chosen base gauging. Use the right-action convention of
\cite{DGP RZ}, Definition 2.7. Put $\zeta=e^{2\pi i/3}$ and
$\eta_b=e^{2\pi i b/9}$, $b\in\{0,1,2\}$, and change the crossed
associator by
$\nu_b(i,j,k)=\zeta^{bk\lfloor(i+j)/3\rfloor}$.
The unit-valued crossed zesting multiplies $\mu_k$ by $\nu_b(i,j,k)$,
$\gamma_{i,j}$ on grade $k$ by its inverse, and leaves crossed
braiding unchanged; see equations (3.19), (3.22)--(3.24) there.
If $X=(T,u)$ has grade $a$, set $u'_h=\eta_b^{ah}u_h$.
The equivariance equation is satisfied precisely because of the
carry in the new compositor. Tensor products acquire the electric
carry $e^{2b\lfloor(i+j)/3\rfloor}$. Thus
\[
\Hom_{\mathcal G_b}(F_b(Y),F_b(X)^n)
\cong\Hom_{\mathcal G_0}
(Y e^{-2b\lfloor na/3\rfloor},X^n).
\]
The braid character changes by $\eta_b^{a^2w}$, and the twist changes
by $\eta_b^{a^2}$. Scalar associator factors cancel on identical
strand grades. These changes agree with the projectors and all
scalar extensions above. Since $H^2(C_3,\Fp^3)=0$ and
$H^3(C_3,\mathbb C^\times)\cong C_3$, these are all allowed choices.

\section*{Acknowledgments}
I thank Cygnus for a complete independent assessment of the mathematical
arguments and their relation to previous work.
The results have not been checked by a human mathematician.

\begin{thebibliography}{99}
\bibitem{AP} J. Achter and R. Pries, The integral monodromy of hyperelliptic
and trielliptic curves, Math. Ann. 338 (2007), 187--206, arXiv:math/0608038.
\bibitem{BJQ} M. Barkeshli, C.-M. Jian and X.-L. Qi, Genons, twist defects,
and projective non-Abelian braiding statistics, arXiv:1208.4834v2.
\bibitem{BS} T. Barmeier and C. Schweigert, A geometric construction for
permutation equivariant categories from modular functors, arXiv:1004.1825.
\bibitem{Atlas} J. Berean-Dutcher, C. C\'ordova, J. Parra-Martinez and
R. Tuddenham, The Anyon Atlas, arXiv:2610.25439.
\bibitem{DM} T. Deshpande and S. Mukhopadhyay, Crossed modular categories
and the Verlinde formula for twisted conformal blocks, arXiv:1909.10799v4.
\bibitem{DGP RZ} C. Delaney, C. Galindo, J. Plavnik, E. Rowell and Q. Zhang,
$G$-crossed braided zesting, J. London Math. Soc. 109 (2024), e12816,
arXiv:2212.05336.
\bibitem{GPR} C. Galindo, J. Plavnik and E. Rowell, Integral
non-group-theoretical modular categories of dimension $p^2q^2$,
Bull. Belg. Math. Soc. Simon Stevin 31 (2024), 516--526, arXiv:2404.03826.
\bibitem{GH} S. Gurevich and R. Hadani, Bounds on certain higher-dimensional
exponential sums via the self-reducibility of the Weil representation,
arXiv:math/0612765v2.
\bibitem{GKRZ} P. Gustafson, A. Kimball, E. Rowell and Q. Zhang,
Braid group representations from twisted tensor products of algebras,
Peking Math. J. 3 (2020), 103--130, arXiv:1906.08153.
\bibitem{GR26} C. Galindo and E. Rowell, Unitary Yang--Baxter operators:
towards a classification, arXiv:2608.16865v1.
\bibitem{Stirling} S. Stirling, Abelian Chern--Simons theory with toral
gauge group, modular tensor categories, and group categories,
arXiv:0807.2857.
\bibitem{Ven} T. N. Venkataramana, Monodromy of cyclic coverings of the
projective line, Invent. Math. 197 (2014), 1--45, arXiv:1204.4778.
\bibitem{ZS} A. E. Zalesskii and V. N. Serezhkin, Linear groups generated
by transvections, Math. USSR-Izv. 10 (1976), 25--46.
\end{thebibliography}
\end{document}
